\section{Benchmarks}\label{sec:benchmarks}

\subsection{Items}

We use four task domains, items listed in Appendix~\ref{app:benchmark_items}:

\begin{itemize}[leftmargin=1.4em]
\item \textbf{poetry\_gen} ($n_{\text{12}}$ items): paired (form, topic, must-avoid) targets across haiku, tanka, sonnet, free verse and a small ghazal slice. POEMetric six-axis scoring~\citep{Ye2025POEMetric}: creativity, lexical diversity, idiosyncrasy, emotional resonance, literary devices, imagery.
\item \textbf{poetry\_interp} ($n_{\text{10}}$ items): a stratified Wittgenstein aspect-shift slice. Each item is a single line of poetry with two pre-registered candidate readings; the model must produce both and the scorer measures aspect-coverage (cosine to each aspect target), aspect-multiplicity (count of aspects $\geq$ 0.40 cosine), novelty (1$-$ retrieval-set similarity).
\item \textbf{aut} ($n_{\text{8}}$ items): the alternative-uses-task component of CreativityPrism~\citep{CreativityPrism2025}. The composite is creativity (cosine novelty against the standard use), lexical-diversity, and feasibility-style heuristics.
\item \textbf{sci\_creativity} ($n_{\text{8}}$ items): scientific-creativity items drawn from BIG-Bench Hard / MacGyverBench style probes~\citep{suzgun2023bbh,Tian2024MacGyver}; each carries 2--3 ``framings'' (analogical, mechanistic, mathematical) the model is asked to produce. The composite weights non-textbook novelty, framing-coverage, and depth-via-length.
\end{itemize}

\subsection{Scoring}

All four scoring functions live in \texttt{benchmarks/scoring.py} and are deterministic: they use \texttt{sentence-transformers/all-MiniLM-L6-v2} (the same embedder PCE uses internally) for cosine-based novelty and aspect-coverage, plus light token-level statistics (lexical diversity, length, presence of imagery cues). We deliberately avoid LLM-as-judge so the contrasts are between observable outputs and not between an evaluator's preference distribution.

\subsection{Why the v0.3 four-arm design controls the v0.2 review's confounds}\label{sec:bench-fairness}

The v0.2 adversarial review identified two confounds the v0.2 four-arm design left unexamined: (a) any v0.2 \texttt{haiku\_cascade} win could simply reflect the \emph{extra compute} purchased by the cascade's K candidates and second pass, and (b) the win could reflect \emph{any} 2-pass revision protocol -- not the specific content of the \iast{vimar\'sa} brief. The v0.3 design replaces v0.2's \texttt{local\_*} arms with two new control arms targeting these confounds:

\begin{itemize}[leftmargin=1.4em]
\item \textbf{\texttt{haiku\_bare\_2K\_scorer}} (extra-compute control). One pass, but with $K' = 2K$ candidates fanned out by \iast{icch\=a} and the same \iast{j\~n\=ana} scorer used by \texttt{haiku\_cascade}. This matches the cascade's compute-budget while denying it the second pass and the \iast{vimar\'sa} brief.
\item \textbf{\texttt{haiku\_generic\_revise\_2pass}} (revision-protocol control). Two passes with \texttt{commit\_policy="always\_revise"}, but the \iast{vimar\'sa}-generated brief is replaced by a fixed generic creative-revise prompt (\texttt{GENERIC\_REVISE\_BRIEF}). Isolates the \emph{content} of the brief from the \emph{existence} of a revision pass.
\end{itemize}

Together the four arms (\texttt{haiku\_bare}, \texttt{haiku\_cascade}, \texttt{haiku\_bare\_2K\_scorer}, \texttt{haiku\_generic\_revise\_2pass}) span: (1) architecture vs.\ nothing (H1.v3--H4.v3), (2) architecture vs.\ \emph{more compute} (H6.v3, the headline fairness test), (3) architecture vs.\ \emph{generic 2-pass} (H7.v3), and (4) within-architecture revision-causality on items the event-gated commit chose to revise (H8.v3). The local-substrate arms are removed from the default matrix per the locked v0.3 scope; the legacy \texttt{local\_*} arms remain callable for backward-compatibility audits.

\subsection{v0.3 prove-gate (single-case validation before pilot)}

Before running the pilot we executed a two-case prove-gate (\texttt{scripts/prove\_gate.py}) on \texttt{tests/fixtures/duck\_rabbit\_textual.json} (Wittgenstein duck-rabbit textual probe; \texttt{id=duck\_rabbit\_textual\_v0.3}) and \texttt{tests/fixtures/aut\_brick.json} (CreativityPrism AUT brick; \texttt{id=aut\_brick\_v0.3}) over both \texttt{haiku\_bare} and \texttt{haiku\_cascade}. Per the v0.3 fixture \texttt{expected\_signals}, the prove-gate now additionally verifies: (i) the per-item \texttt{IntegrityProbe} passes (no Claude-Code framing leaks into the inner CLI subprocess), (ii) every Haiku surface and revision text passes \texttt{LEAKAGE\_REGEX} (negation-aware), (iii) the duck-rabbit fixture forces \iast{vimar\'sa}-event commit to revision with $|\Delta F| \geq 0.01$ (non-degenerate aspect-conditioned BMR), (iv) the AUT brick (no aspects) commits the draft with $\Delta F = 0$ (event-gated commit semantics), and (v) when \texttt{committed=="revision"} the revision text differs from the draft text. Local-substrate arms are not exercised in the v0.3 prove-gate.
